\chapter*{Annexes}
\addcontentsline{toc}{chapter}{Annexes}
\markboth{Annexes}{Annexes}

\section*{Annexe A : Guide d'installation et configuration}

Pour déployer ManageraHub en local pour le développement, veuillez suivre ces étapes :

\subsection*{1. Prérequis système}
\begin{itemize}
    \item Python 3.9 installé.
    \item Gestionnaire de paquets \texttt{pip} et utilitaire de virtualisation \texttt{venv}.
    \item Git installé pour le clonage du dépôt.
\end{itemize}

\subsection*{2. Étapes de déploiement}
\begin{lstlisting}[language=bash, caption=Commandes de configuration du projet]
# 1. Cloner le depot du projet
git clone https://github.com/TchichAdam/managerahub-django.git
cd managerahub-django

# 2. Creer et activer l'environnement virtuel
python -m venv venv
# Sur Windows :
venv\Scripts\activate
# Sur Linux/macOS :
source venv/bin/activate

# 3. Installer les dependances
pip install -r requirements.txt

# 4. Creer les fichiers de migration de base de donnees
python manage.py makemigrations app1

# 5. Appliquer les migrations
python manage.py migrate

# 6. Creer un compte administrateur principal
python manage.py createsuperuser

# 7. Lancer le serveur de developpement local
python manage.py runserver
\end{lstlisting}

Une fois le serveur démarré, l'application est accessible à l'adresse suivante dans n'importe quel navigateur web : \url{http://127.0.0.1:8000/}.

\subsection*{3. Comptes de test prédéfinis}
\begin{itemize}
    \item \textbf{Compte Administrateur :} \texttt{admin@managerahub.com} / Mot de passe : \texttt{admin123}
    \item \textbf{Candidat de test :} \texttt{candidate@}\allowbreak\texttt{test.com} / Passe : \texttt{candidate123}
    \item \textbf{Compte Entreprise de test (en attente) :} \texttt{entreprise@test.com} / Mot de passe : \texttt{entreprise123}
\end{itemize}
\textit{Note de sécurité importante : Ces comptes de test et leurs mots de passe associés sont réservés exclusivement à l'environnement de développement local (dev-only) et doivent impérativement être supprimés ou modifiés en production.}

\newpage

\section*{Annexe B : Structure des gabarits HTML (Templates)}
L'architecture de présentation (frontend) de ManageraHub est conçue selon un modèle d'héritage de templates Django, favorisant la réutilisabilité du code et le respect du principe DRY (Don't Repeat Yourself).

\subsection*{1. Fichier de base global (templates/base.html)}
Ce fichier contient la structure HTML5 globale commune, l'intégration des polices Cormorant Garamond et Inter via Google Fonts, et le point d'inclusion de la feuille de style principale.

\begin{lstlisting}[language=HTML, caption=Structure réelle du template base.html]
{% load static %}
<!DOCTYPE html>
<html lang="en">
<head>
    <meta charset="UTF-8">
    <meta name="viewport" content="width=device-width, initial-scale=1.0">
    <title>{% block title %}ManageraHub{% endblock %}</title>
    <meta name="description" content="ManageraHub is a modern recruitment platform that connects candidates, companies, and administrators in one professional experience.">
    <link rel="preconnect" href="https://fonts.googleapis.com">
    <link rel="preconnect" href="https://fonts.gstatic.com" crossorigin>
    <link href="https://fonts.googleapis.com/css2?family=Cormorant+Garamond:wght@400;500;600;700&family=Inter:wght@400;500;600;700;800&display=swap" rel="stylesheet">
    <link rel="stylesheet" href="{% static 'css/main.css' %}?v=20260615-feed-v11">
</head>
<body>
    {% block content %}{% endblock %}
    <script src="{% static 'js/main.js' %}?v=20260520-2"></script>
    {% block extra_scripts %}{% endblock %}
</body>
</html>
\end{lstlisting}

\subsection*{2. Fichier de tableau de bord Candidat (\texttt{candidate/}\allowbreak\texttt{dashboard.html})}
Ce gabarit affiche le profil du candidat, la complétion de son profil et ses candidatures récentes. Il hérite du gabarit \texttt{candidate/}\allowbreak\texttt{base.html}.

\begin{lstlisting}[language=HTML, caption=Template du Dashboard Candidat (Extrait)]
{% extends "candidate/base.html" %}

{% block title %}Candidate Dashboard | ManageraHub{% endblock %}

{% block candidate_main %}
<section class="candidate-hero-card candidate-welcome">
    <div class="candidate-welcome-text">
        <p class="candidate-card-kicker">Candidate Dashboard</p>
        <h1>Welcome back, {{ candidate_display_name }}.</h1>
        <p>Track your applications, finish your profile, and discover fresh opportunities.</p>
        <div class="candidate-hero-actions">
            <a href="{% url 'candidate_profile' %}" class="pill-button pill-button-dark">Complete my profile</a>
            <a href="{% url 'candidate_job_offers' %}" class="pill-button pill-button-light">Browse job offers</a>
        </div>
    </div>
    <div class="candidate-welcome-aside">
        <div class="candidate-progress-ring" style="--p: {{ candidate_completion_percent }};">
            <div class="candidate-progress-ring-center">
                <strong>{{ candidate_completion_percent }}%</strong>
                <span>Profile</span>
            </div>
        </div>
    </div>
</section>

<section class="candidate-dash-grid">
    <div class="candidate-dash-main">
        <article class="candidate-panel-card">
            <p class="candidate-card-kicker">Recent applications</p>
            {% if applications %}
            <div class="candidate-mini-stack">
                {% for application in applications %}
                <div class="candidate-mini-entry">
                    <strong>{{ application.job_offer.title }}</strong>
                    <span>{{ application.job_offer.company_name }} - {{ application.get_status_display }}</span>
                </div>
                {% endfor %}
            </div>
            {% else %}
            <p>You have not applied to a job yet. Start from the Job Offers page.</p>
            {% endif %}
        </article>
    </div>
</section>
{% endblock %}
\end{lstlisting}

\newpage

\section*{Annexe C : Validations de formulaires et logique de sécurité}
Pour garantir la cohérence des données, notamment pour le marché marocain, des mécanismes de validation stricts ont été mis en œuvre.

\subsection*{1. Validateur personnalisé d'ICE marocain (app1/forms.py)}
L'ICE (Identifiant Commun de l'Entreprise) doit être composé de 15 chiffres. Voici comment ce contrôle est géré dans notre formulaire de modification du profil de l'entreprise :

\begin{lstlisting}[language=Python, caption=Validation de l'ICE dans CompanyProfileForm]
from django import forms
from .models import CompanyProfile

class CompanyProfileForm(forms.ModelForm):
    class Meta:
        model = CompanyProfile
        fields = [
            "company_name",
            "industry",
            "company_size",
            "phone_number",
            "city",
            "country",
            "website",
            "description",
            "logo",
            "background_image",
            "ice",
            "rc_number",
            "legal_document",
        ]
        
    def clean_ice(self):
        ice = self.cleaned_data.get("ice")
        if ice:
            ice = ice.strip()
            if not ice.isdigit() or len(ice) != 15:
                raise forms.ValidationError("The ICE number must contain exactly 15 digits.")
        return ice
\end{lstlisting}
\newpage
